% Options for packages loaded elsewhere
\PassOptionsToPackage{unicode}{hyperref}
\PassOptionsToPackage{hyphens}{url}
%
\documentclass[
]{article}
\usepackage{amsmath,amssymb}
\usepackage{iftex}
\ifPDFTeX
  \usepackage[T1]{fontenc}
  \usepackage[utf8]{inputenc}
  \usepackage{textcomp} % provide euro and other symbols
\else % if luatex or xetex
  \usepackage{unicode-math} % this also loads fontspec
  \defaultfontfeatures{Scale=MatchLowercase}
  \defaultfontfeatures[\rmfamily]{Ligatures=TeX,Scale=1}
\fi
\usepackage{lmodern}
\ifPDFTeX\else
  % xetex/luatex font selection
\fi
% Use upquote if available, for straight quotes in verbatim environments
\IfFileExists{upquote.sty}{\usepackage{upquote}}{}
\IfFileExists{microtype.sty}{% use microtype if available
  \usepackage[]{microtype}
  \UseMicrotypeSet[protrusion]{basicmath} % disable protrusion for tt fonts
}{}
\makeatletter
\@ifundefined{KOMAClassName}{% if non-KOMA class
  \IfFileExists{parskip.sty}{%
    \usepackage{parskip}
  }{% else
    \setlength{\parindent}{0pt}
    \setlength{\parskip}{6pt plus 2pt minus 1pt}}
}{% if KOMA class
  \KOMAoptions{parskip=half}}
\makeatother
\usepackage{xcolor}
\usepackage[margin=1in]{geometry}
\usepackage{color}
\usepackage{fancyvrb}
\newcommand{\VerbBar}{|}
\newcommand{\VERB}{\Verb[commandchars=\\\{\}]}
\DefineVerbatimEnvironment{Highlighting}{Verbatim}{commandchars=\\\{\}}
% Add ',fontsize=\small' for more characters per line
\usepackage{framed}
\definecolor{shadecolor}{RGB}{248,248,248}
\newenvironment{Shaded}{\begin{snugshade}}{\end{snugshade}}
\newcommand{\AlertTok}[1]{\textcolor[rgb]{0.94,0.16,0.16}{#1}}
\newcommand{\AnnotationTok}[1]{\textcolor[rgb]{0.56,0.35,0.01}{\textbf{\textit{#1}}}}
\newcommand{\AttributeTok}[1]{\textcolor[rgb]{0.13,0.29,0.53}{#1}}
\newcommand{\BaseNTok}[1]{\textcolor[rgb]{0.00,0.00,0.81}{#1}}
\newcommand{\BuiltInTok}[1]{#1}
\newcommand{\CharTok}[1]{\textcolor[rgb]{0.31,0.60,0.02}{#1}}
\newcommand{\CommentTok}[1]{\textcolor[rgb]{0.56,0.35,0.01}{\textit{#1}}}
\newcommand{\CommentVarTok}[1]{\textcolor[rgb]{0.56,0.35,0.01}{\textbf{\textit{#1}}}}
\newcommand{\ConstantTok}[1]{\textcolor[rgb]{0.56,0.35,0.01}{#1}}
\newcommand{\ControlFlowTok}[1]{\textcolor[rgb]{0.13,0.29,0.53}{\textbf{#1}}}
\newcommand{\DataTypeTok}[1]{\textcolor[rgb]{0.13,0.29,0.53}{#1}}
\newcommand{\DecValTok}[1]{\textcolor[rgb]{0.00,0.00,0.81}{#1}}
\newcommand{\DocumentationTok}[1]{\textcolor[rgb]{0.56,0.35,0.01}{\textbf{\textit{#1}}}}
\newcommand{\ErrorTok}[1]{\textcolor[rgb]{0.64,0.00,0.00}{\textbf{#1}}}
\newcommand{\ExtensionTok}[1]{#1}
\newcommand{\FloatTok}[1]{\textcolor[rgb]{0.00,0.00,0.81}{#1}}
\newcommand{\FunctionTok}[1]{\textcolor[rgb]{0.13,0.29,0.53}{\textbf{#1}}}
\newcommand{\ImportTok}[1]{#1}
\newcommand{\InformationTok}[1]{\textcolor[rgb]{0.56,0.35,0.01}{\textbf{\textit{#1}}}}
\newcommand{\KeywordTok}[1]{\textcolor[rgb]{0.13,0.29,0.53}{\textbf{#1}}}
\newcommand{\NormalTok}[1]{#1}
\newcommand{\OperatorTok}[1]{\textcolor[rgb]{0.81,0.36,0.00}{\textbf{#1}}}
\newcommand{\OtherTok}[1]{\textcolor[rgb]{0.56,0.35,0.01}{#1}}
\newcommand{\PreprocessorTok}[1]{\textcolor[rgb]{0.56,0.35,0.01}{\textit{#1}}}
\newcommand{\RegionMarkerTok}[1]{#1}
\newcommand{\SpecialCharTok}[1]{\textcolor[rgb]{0.81,0.36,0.00}{\textbf{#1}}}
\newcommand{\SpecialStringTok}[1]{\textcolor[rgb]{0.31,0.60,0.02}{#1}}
\newcommand{\StringTok}[1]{\textcolor[rgb]{0.31,0.60,0.02}{#1}}
\newcommand{\VariableTok}[1]{\textcolor[rgb]{0.00,0.00,0.00}{#1}}
\newcommand{\VerbatimStringTok}[1]{\textcolor[rgb]{0.31,0.60,0.02}{#1}}
\newcommand{\WarningTok}[1]{\textcolor[rgb]{0.56,0.35,0.01}{\textbf{\textit{#1}}}}
\usepackage{graphicx}
\makeatletter
\def\maxwidth{\ifdim\Gin@nat@width>\linewidth\linewidth\else\Gin@nat@width\fi}
\def\maxheight{\ifdim\Gin@nat@height>\textheight\textheight\else\Gin@nat@height\fi}
\makeatother
% Scale images if necessary, so that they will not overflow the page
% margins by default, and it is still possible to overwrite the defaults
% using explicit options in \includegraphics[width, height, ...]{}
\setkeys{Gin}{width=\maxwidth,height=\maxheight,keepaspectratio}
% Set default figure placement to htbp
\makeatletter
\def\fps@figure{htbp}
\makeatother
\setlength{\emergencystretch}{3em} % prevent overfull lines
\providecommand{\tightlist}{%
  \setlength{\itemsep}{0pt}\setlength{\parskip}{0pt}}
\setcounter{secnumdepth}{-\maxdimen} % remove section numbering
\ifLuaTeX
  \usepackage{selnolig}  % disable illegal ligatures
\fi
\IfFileExists{bookmark.sty}{\usepackage{bookmark}}{\usepackage{hyperref}}
\IfFileExists{xurl.sty}{\usepackage{xurl}}{} % add URL line breaks if available
\urlstyle{same}
\hypersetup{
  pdftitle={STAT 118: Notes B},
  hidelinks,
  pdfcreator={LaTeX via pandoc}}

\title{STAT 118: Notes B}
\usepackage{etoolbox}
\makeatletter
\providecommand{\subtitle}[1]{% add subtitle to \maketitle
  \apptocmd{\@title}{\par {\large #1 \par}}{}{}
}
\makeatother
\subtitle{Wrangling data with dplyr: filter, select, arrange}
\author{}
\date{\vspace{-2.5em}}

\begin{document}
\maketitle

\hypertarget{importing-data}{%
\section{Importing Data}\label{importing-data}}

In this class, we are going to be working with a dataset relating to the
languages spoken at home by Canadian Residents. Many Indigenous peoples
exist in Canada with their own languages and cultures. Sadly,
colonization has led to the loss of many of these languages. This data
is a subset of data collected during the 2016 census.

What is a .csv file?

How do we import it into R?

\begin{Shaded}
\begin{Highlighting}[]
\CommentTok{\#can\_lang.csv needs to be saved in the same directory as this file. }
\CommentTok{\#can\_lang \textless{}{-} read.csv("can\_lang.csv")}
\end{Highlighting}
\end{Shaded}

Alternatively, you can download it directly from the internet. Github
user \texttt{ttimbers} hosts this file to share with the public.

\begin{Shaded}
\begin{Highlighting}[]
\NormalTok{can\_lang }\OtherTok{\textless{}{-}} \FunctionTok{read.csv}\NormalTok{(}\StringTok{"https://raw.githubusercontent.com/ttimbers/canlang/master/inst/extdata/can\_lang.csv"}\NormalTok{)}
\end{Highlighting}
\end{Shaded}

Let's take a look at this data for a minute to see what information has
been recorded.

\hypertarget{installing-and-using-packages}{%
\section{Installing and Using
Packages}\label{installing-and-using-packages}}

Sometimes everything we need (data, functions, etc) are not available in
base R. In R, expert users will package up useful things like data and
functions into packages that be download and used.

First, you need to download the package from the right hand menu
--\textgreater{} You only need to do this once.

In each new .Rmd document, you need to call any packages you want to use
but adding the code \texttt{library(packagename)} inside an R chunk.

For example, in this class we will use the \texttt{tidyverse} package a
lot.

\begin{Shaded}
\begin{Highlighting}[]
\FunctionTok{library}\NormalTok{(tidyverse)}
\end{Highlighting}
\end{Shaded}

\hypertarget{dplyr}{%
\section{\texorpdfstring{\texttt{dplyr}}{dplyr}}\label{dplyr}}

There are actually many commonly used packages wrapped up inside one
\texttt{tidyverse} package.

Today we are specifically going to be talking about the package
\texttt{dplyr} which is useful to manipulating data sets.

\hypertarget{filter}{%
\section{\texorpdfstring{\texttt{filter}}{filter}}\label{filter}}

We can use the \texttt{filter} function to extract \textbf{\emph{rows}}
from the data that have a particular characteristic.

For example, we may be interested in only looking at only the languages
in this dataset that are Aboriginal languages.

\begin{Shaded}
\begin{Highlighting}[]
\CommentTok{\#start with the can\_lang dataset, the pipe "\%" means apply the action on the following line to the previous line. In this case, pick out only the rows were the category variable is "Aboriginal languages" }




\DocumentationTok{\#\#note the aboriginal languages is text/categorical and so quotation marks are needed. }
\DocumentationTok{\#\#R doesn\textquotesingle{}t care about whether they are double quotation marks (") or single (\textquotesingle{}). They work the same. }
\CommentTok{\# If we don\textquotesingle{}t assign it to an object, then it just prints out for us to see! }
\end{Highlighting}
\end{Shaded}

\begin{Shaded}
\begin{Highlighting}[]
\CommentTok{\#oftentimes, we want to take our subset and give it a new name. This takes our subset and assigns it to a new dataset called \textasciigrave{}aboriginal\_lang\textasciigrave{}. }




\CommentTok{\#Notice if you assign it to an object that it doesn\textquotesingle{}t print out the contents. }
\CommentTok{\# You\textquotesingle{}ll see the new object in your environment on the top right {-}{-}{-}\textgreater{} }
\CommentTok{\# If you click on the word \textasciigrave{}aboriginal languages\textasciigrave{} (not the blue play button) it will open the object so you can see what is saved inside. }
\end{Highlighting}
\end{Shaded}

It can also be used with numeric criteria.

Suppose we want a list of all the languages in Canada that are spoken by
less than 100 people as their mother tongue.

The logical operators are given below:

\includegraphics[width=0.5\textwidth,height=\textheight]{https://ehsanx.github.io/intro2R/images/logical.png}

\hypertarget{select}{%
\section{\texorpdfstring{\texttt{select}}{select}}\label{select}}

select is used to extract only certain \textbf{\emph{columns}}. For
example, perhaps we only want to print out a list names of the
aboriginal languages (language column).

We can combine criteria together as well in one command with multiple
pipes:

\hypertarget{arrange}{%
\section{\texorpdfstring{\texttt{arrange}}{arrange}}\label{arrange}}

The \texttt{arrange} function allows us to order the rows of the data
frame by the values of a particular column.

For example, arrange all the aboriginal languages in canada by from most
to least spoken as mother tongue.

\hypertarget{slice}{%
\section{\texorpdfstring{\texttt{slice}}{slice}}\label{slice}}

The slice function will allow us to pick only a subset of the rows based
on their numeric order (1st through last).

For example, if I want a list of the 10 most commonly spoken aboriginal
languages.

\hypertarget{badges-earned}{%
\paragraph[Badges Earned ]{\texorpdfstring{Badges Earned
\protect\includegraphics[width=1in,height=\textheight]{https://github.com/rstudio/hex-stickers/raw/main/thumbs/pipe.png}
\protect\includegraphics[width=1in,height=\textheight]{https://github.com/rstudio/hex-stickers/raw/main/thumbs/dplyr.png}}{Badges Earned  }}\label{badges-earned}}

\hypertarget{brain-break}{%
\section{Brain Break}\label{brain-break}}

\begin{itemize}
\item
  \href{https://www.youtube.com/watch?v=99-LoEkAA3w}{Students at Allison
  Bernard Memorial High School in Eskasoni, Cape Breton recorded Paul
  McCartney's Blackbird in their native Mi'kmaq language.}
\item
  \href{https://www.youtube.com/watch?v=wW0gpo2deKg}{The Jerry Cans are
  a band from Iqaluit, Nunavut who combine traditional Inuit throat
  singing with folk music and country rock. Their music is largely
  written in Inuktitut (the indigineous language of the Inuit)}
\end{itemize}

\end{document}
